% !TeX root = ../../ONMIPA Alin.tex
% Acuan: Zhang (2009), Bab 3, definisi dan fakta serta Soal 3.79--3.112;
% Zhang (2011), bagian 1.3. Contoh dan pembahasan disusun untuk handout ini.
\chapter{Transformasi Linear}
\label{ch:transformasi-linear}

Transformasi linear mempertahankan kombinasi linear. Karena itu,
perilakunya pada sebuah basis menentukan perilakunya di seluruh
ruang. Bab ini menghubungkan struktur ruang vektor dengan perhitungan
matriks, lalu memakai hubungan tersebut untuk menganalisis kernel,
citra, proyeksi, serta pangkat suatu operator.
Semua ruang dan transformasi menggunakan medan skalar yang sama,
yaitu $\F=\R$ atau $\C$, kecuali jika dinyatakan lain.

\section{Definisi dan Penentuan pada Basis}

\begin{definisi}[][Transformasi linear][trans:def]
    Pemetaan $T:V\to W$ disebut \emph{transformasi linear} jika
    untuk semua $u,v\in V$ dan $a,b\in\F$ berlaku
    \[
        T(au+bv)=aT(u)+bT(v).
    \]
    Transformasi linear dari $V$ ke dirinya sendiri disebut
    \emph{operator linear} pada $V$.
\end{definisi}

Definisi tersebut ekuivalen dengan aditivitas
$T(u+v)=T(u)+T(v)$ dan homogenitas $T(av)=aT(v)$.
Dengan mengambil koefisien nol, diperoleh $T(0_V)=0_W$.
Induksi kemudian memberi
\[
    T\left(\sum_{j=1}^r a_jv_j\right)
      =\sum_{j=1}^r a_jT(v_j).
\]
Jadi setiap hubungan linear antarvektor harus tetap berlaku
setelah transformasi diterapkan. Syarat $T(0)=0$ saja belum
cukup: pemetaan $T(x)=x^2$ pada $\R$ memenuhi syarat itu,
tetapi gagal pada penjumlahan.

\begin{exam}[Linear terhadap medan yang mana?]
    Pada $\C$, tinjau $T(z)=\overline z$ dan $S(z)=\operatorname{Re}z$,
    dengan nilai $S$ juga dipandang sebagai bilangan kompleks.
    Apakah keduanya linear atas $\R$? Atas $\C$?
\end{exam}
\begin{solution}
    Konjugasi dan bagian real mempertahankan penjumlahan dan
    perkalian oleh skalar real, sehingga keduanya linear atas $\R$.
    Namun, $T(i)=-i\ne iT(1)$ dan $S(i)=0\ne iS(1)$.
    Jadi keduanya tidak linear atas $\C$.
    Pilihan medan skalar merupakan bagian dari pernyataan linearitas.
\end{solution}

Pemetaan $x\mapsto Ax+b$ dari $\F^n$ ke $\F^m$ linear tepat
ketika $b=0$; syarat ini langsung diperoleh dari nilai di nol.
Pemetaan dengan $b\ne0$ disebut pemetaan afin.
Contoh linear lain mencakup evaluasi $p\mapsto p(a)$,
turunan $p\mapsto p'$, serta pemetaan matriks $X\mapsto AXB$
untuk $A,B$ tetap dengan ukuran yang sesuai.

\begin{teorema}[][Menentukan transformasi melalui basis][trans:basis]
    Misalkan $\mathcal B=(v_1,\ldots,v_n)$ basis bagi $V$.
    Untuk sembarang $w_1,\ldots,w_n\in W$, terdapat tepat satu
    transformasi linear $T:V\to W$ dengan $T(v_j)=w_j$.
    Transformasi tersebut diberikan oleh
    \[
        T\left(\sum_{j=1}^n a_jv_j\right)=\sum_{j=1}^n a_jw_j.
    \]
\end{teorema}
\begin{proof}
    Koordinat terhadap basis tunggal, sehingga rumus tersebut
    memberikan satu nilai yang tidak bergantung pada pilihan
    representasi. Rumusnya mempertahankan kombinasi linear.
    Setiap transformasi linear dengan nilai basis yang ditentukan
    harus memenuhi rumus itu, sehingga transformasinya unik.
\end{proof}

Jika nilai transformasi ditentukan pada daftar yang hanya
membangun ruang, semua hubungan pada daftar itu harus dihormati.
Secara tepat, aturan $v_j\mapsto w_j$ memperluas menjadi
transformasi pada $\operatorname{span}\{v_j\}$ jika dan hanya jika
\[
    \sum a_jv_j=0\quad\Longrightarrow\quad\sum a_jw_j=0.
\]
Syarat ini perlu karena linearitas, dan cukup karena dua
representasi dari vektor yang sama mempunyai selisih nol.

\begin{exam}[Data pada tiga vektor yang bukan basis standar]
    Carilah $T:\R^3\to\R^2$ yang memenuhi
    \[
        T(e_1+e_2)=(1,0),\quad
        T(e_2+e_3)=(0,1),\quad
        T(e_3+e_1)=(1,1).
    \]
\end{exam}
\begin{solution}
    Ketiga vektor sumber bebas linear, sebab determinan matriks
    kolomnya adalah $2$. Jadi transformasi ada dan unik.
    Tulis $a=T(e_1)$, $b=T(e_2)$, $c=T(e_3)$.
    Sistem $a+b=(1,0)$, $b+c=(0,1)$, $c+a=(1,1)$
    memberi $a=(1,0)$, $b=(0,0)$, $c=(0,1)$.
    Maka $T(x,y,z)=(x,z)$.
\end{solution}

\section{Kernel, Citra, dan Teorema Rank--Nulitas}

\begin{definisi}[][Kernel, citra, rank, dan nulitas][trans:kernel]
    Untuk transformasi linear $T:V\to W$, definisikan
    \[
        \ker T=\{v\in V:T(v)=0\},\qquad
        \operatorname{im}T=\{T(v):v\in V\}.
    \]
    Jika dimensinya hingga, tulis
    $\nul T=\dim\ker T$ dan $\rank T=\dim\operatorname{im}T$.
\end{definisi}

Kernel merupakan subruang \emph{domain} $V$, sedangkan citra
merupakan subruang \emph{kodomain} $W$.
Untuk kernel, $T(au+bv)=0$ jika $T(u)=T(v)=0$.
Untuk citra, kombinasi $aT(u)+bT(v)$ sama dengan $T(au+bv)$.
Kedua himpunan memuat nol, sehingga keduanya memenuhi uji subruang.

\begin{teorema}[][Teorema rank--nulitas][trans:rank-nulitas]
    Jika $T:V\to W$ linear dan $\dim V=n<\infty$, maka
    \[
        \rank T+\nul T=n.
    \]
    Lebih tepatnya, ambil basis $(u_1,\ldots,u_k)$ bagi $\ker T$,
    lalu perluas menjadi basis
    \[
        (u_1,\ldots,u_k,v_1,\ldots,v_{n-k})
    \]
    bagi $V$. Daftar $(T(v_1),\ldots,T(v_{n-k}))$ merupakan basis citra.
\end{teorema}
\begin{proof}
    Ekspansi vektor domain dalam basis yang diberikan menunjukkan
    bahwa citra dibangun oleh $T(v_1),\ldots,T(v_{n-k})$.
    Jika $\sum a_jT(v_j)=0$, maka $\sum a_jv_j\in\ker T$
    dan dapat dituliskan sebagai kombinasi $u_1,\ldots,u_k$.
    Kebebasan basis $V$ memaksa semua $a_j=0$.
    Maka dimensi citra adalah $n-k$, sedangkan nulitasnya $k$.
\end{proof}

\begin{prop}[Injektif, surjektif, dan kebebasan citra]
    Transformasi $T$ injektif jika dan hanya jika $\ker T=\{0\}$,
    dan surjektif jika dan hanya jika $\operatorname{im}T=W$.
    Jika $T(v_1),\ldots,T(v_r)$ bebas linear, maka $v_1,\ldots,v_r$
    bebas linear. Arah sebaliknya berlaku jika $T$ injektif.
\end{prop}
\begin{proof}
    Persamaan $T(u)=T(v)$ ekuivalen dengan $u-v\in\ker T$,
    sehingga kriteria injektivitas mengikuti langsung.
    Kriteria surjektivitas adalah definisi citra.
    Untuk kebebasan, terapkan $T$ pada suatu hubungan linear
    vektor sumber. Sebaliknya, jika citra suatu kombinasi nol
    dan $T$ injektif, kombinasi sumbernya juga nol.
\end{proof}

Jika $\dim V=\dim W=n<\infty$, teorema rank--nulitas memberi
ekuivalensi: $T$ injektif, $T$ surjektif, dan $T$ bijektif.
Untuk dimensi berbeda, injektivitas hanya mengharuskan
$\dim V\le\dim W$, sedangkan surjektivitas mengharuskan
$\dim V\ge\dim W$. Ketaksamaan dimensi saja tidak menjamin
sifat tersebut bagi transformasi yang diberikan.

\begin{exam}[Evaluasi polinom pada beberapa titik]
    Untuk titik berbeda $a_1,\ldots,a_m\in\F$, $m\ge1$, tinjau
    \[
        E:\mathcal P_{\le n}(\F)\to\F^m,\qquad
        E(p)=(p(a_1),\ldots,p(a_m))^T.
    \]
    Tentukan rank, nulitas, dan kapan $E$ injektif atau surjektif.
\end{exam}
\begin{solution}
    Kernel terdiri atas polinom yang habis dibagi
    $q(t)=\prod_{j=1}^m(t-a_j)$.
    Jika $m\le n$, kernel adalah
    $q\mathcal P_{\le n-m}(\F)$, berdimensi $n-m+1$.
    Jika $m>n$, polinom berderajat paling besar $n$ dengan
    $m$ akar berbeda harus nol. Jadi
    \[
        \nul E=\max\{n+1-m,0\},\qquad
        \rank E=\min\{n+1,m\}.
    \]
    Transformasi injektif tepat ketika $m\ge n+1$ dan surjektif
    tepat ketika $m\le n+1$. Untuk $m=n+1$, interpolasi
    Lagrange memberikan inversnya secara eksplisit.
\end{solution}

\subsection{Serat solusi dan prapeta subruang}
Jika $T(v_0)=w$, seluruh solusi $T(v)=w$ adalah
\[
    T^{-1}(\{w\})=v_0+\ker T.
\]
Memang, $T(v)=T(v_0)$ jika dan hanya jika $v-v_0\in\ker T$.
Jadi suatu persamaan linear konsisten tepat ketika
$w\in\operatorname{im}T$, dan jika konsisten, solusinya unik
tepat ketika kernel hanya nol.

Untuk subruang $Y\subseteq W$, prapeta
$T^{-1}(Y)=\{v:T(v)\in Y\}$ selalu subruang dari $V$.
Notasi prapeta ini tidak mensyaratkan bahwa $T$ mempunyai invers.
Pembatasan $T$ pada $T^{-1}(Y)$ mempunyai kernel $\ker T$
dan citra $Y\cap\operatorname{im}T$. Maka, jika $V$ berdimensi hingga,
\[
    \dim T^{-1}(Y)=\nul T+\dim(Y\cap\operatorname{im}T).
\]
Demikian pula, untuk subruang $U\subseteq V$,
$\dim T(U)=\dim U-\dim(U\cap\ker T)$.

\begin{exam}[Jumlah dimensi tidak menjamin jumlah langsung]
    Pada $\R^2$, definisikan $T(x,y)=(y,0)$.
    Tentukan kernel, citra, dan jumlah keduanya.
\end{exam}
\begin{solution}
    Kernel dan citra keduanya $\operatorname{span}\{e_1\}$.
    Walaupun $\nul T+\rank T=1+1=2$, jumlah kedua subruang
    hanya garis $\operatorname{span}\{e_1\}$, bukan seluruh $\R^2$.
    Teorema rank--nulitas menyatakan kesamaan bilangan dimensi;
    ia tidak menyatakan bahwa domain adalah jumlah langsung
    kernel dan citra. Untuk domain dan kodomain berbeda,
    penjumlahan kernel dan citra bahkan belum tentu bermakna.
\end{solution}

\section{Representasi Matriks dan Pergantian Basis}

\begin{definisi}[][Matriks suatu transformasi][trans:matriks]
    Misalkan $\mathcal B=(b_1,\ldots,b_n)$ basis $V$ dan
    $\mathcal C=(c_1,\ldots,c_m)$ basis $W$.
    Matriks transformasi $T:V\to W$ relatif terhadap kedua basis adalah
    \[
        [T]_{\mathcal C\leftarrow\mathcal B}
        =\bigl[[T(b_1)]_{\mathcal C}\ \cdots\
        [T(b_n)]_{\mathcal C}\bigr]\in\Mat_{m\times n}(\F).
    \]
\end{definisi}

Dengan menuliskan $v=\sum x_jb_j$, linearitas memberi
\[
    [T(v)]_{\mathcal C}
      =[T]_{\mathcal C\leftarrow\mathcal B}[v]_{\mathcal B}.
\]
Jadi kolom ke-$j$ berisi \emph{koordinat citra vektor basis
ke-$j$}, bukan barisnya. Untuk operator dengan basis domain dan
kodomain sama, notasi disingkat menjadi $[T]_{\mathcal B}$.
Rank dan nulitas matriks ini sama dengan rank dan nulitas $T$,
karena persamaan koordinat mempertahankan hubungan linear.

\begin{exam}[Turunan dalam dua basis polinom]
    Pada $V=\mathcal P_{\le3}(\R)$, letakkan $D(p)=p'$.
    Tentukan matriks $D$ dalam basis
    $\mathcal B=(1,t,t^2,t^3)$ dan
    $\mathcal C=(1,t,t^2/2,t^3/6)$.
\end{exam}
\begin{solution}
    Turunan anggota basis langsung memberi
    \[
        [D]_{\mathcal B}
        =\begin{pmatrix}0&1&0&0\\0&0&2&0\\0&0&0&3\\0&0&0&0\end{pmatrix},
        \qquad
        [D]_{\mathcal C}
        =\begin{pmatrix}0&1&0&0\\0&0&1&0\\0&0&0&1\\0&0&0&0\end{pmatrix}.
    \]
    Kedua matriks menyatakan operator yang sama.
    Basis kedua menyerap faktor dari penurunan pangkat.
    Kernel dibangun oleh $1$, citra adalah $\mathcal P_{\le2}(\R)$,
    dan $D^4=0$ pada ruang ini.
\end{solution}

\begin{teorema}[][Rumus pergantian basis][trans:ganti-basis]
    Tulis $A=[T]_{\mathcal C\leftarrow\mathcal B}$,
    $P=P_{\mathcal B\leftarrow\mathcal B'}$, dan
    $Q=P_{\mathcal C\leftarrow\mathcal C'}$.
    Maka
    \[
        [T]_{\mathcal C'\leftarrow\mathcal B'}=Q^{-1}AP.
    \]
    Khusus untuk operator dengan perubahan basis yang sama
    pada domain dan kodomain, matriks baru adalah $P^{-1}AP$.
\end{teorema}
\begin{proof}
    Untuk setiap $v$, lakukan perubahan koordinat berurutan:
    \[
        [T(v)]_{\mathcal C'}
        =Q^{-1}[T(v)]_{\mathcal C}
        =Q^{-1}A[v]_{\mathcal B}
        =Q^{-1}AP[v]_{\mathcal B'}.
    \]
    Keunikan matriks representasi memberikan rumus yang dinyatakan.
\end{proof}

Rumus ini menjelaskan mengapa keserupaan matriks muncul pada
operator linear: dua matriks serupa merepresentasikan operator
yang sama dalam basis berbeda. Akibatnya, determinan, trace,
polinom karakteristik, dan nilai eigen operator tidak bergantung
pada pilihan basis. Untuk transformasi antar dua ruang dengan
perubahan basis terpisah, bentuk umumnya tetap $Q^{-1}AP$;
bentuk itu tidak harus merupakan keserupaan.

\section{Komposisi, Invers, dan Ruang Transformasi}

Jika $T:V\to W$ dan $S:W\to Z$ linear, komposisi
$ST=S\circ T$ linear dan berarti \emph{menerapkan $T$ lebih dahulu}.
Dengan basis yang sesuai,
\[
    [ST]_{\mathcal D\leftarrow\mathcal B}
      =[S]_{\mathcal D\leftarrow\mathcal C}
       [T]_{\mathcal C\leftarrow\mathcal B}.
\]
Urutan perkalian mengikuti urutan komposisi; dua operator pada
ruang yang sama tidak harus komutatif.

\begin{teorema}[][Rank komposisi][trans:komposisi]
    Untuk $T:V\to W$ dan $S:W\to Z$ dengan semua ruang
    berdimensi hingga,
    \[
        \rank(ST)=\rank T-\dim(\operatorname{im}T\cap\ker S).
    \]
    Akibatnya,
    \[
        \rank S+\rank T-\dim W
        \le\rank(ST)\le\min\{\rank S,\rank T\}.
    \]
\end{teorema}
\begin{proof}
    Batasi $S$ pada $\operatorname{im}T$.
    Kernel pembatasannya adalah $\operatorname{im}T\cap\ker S$,
    sedangkan citranya $\operatorname{im}(ST)$.
    Teorema rank--nulitas memberikan rumus utama.
    Karena dimensi irisan paling besar $\dim\ker S
    =\dim W-\rank S$, diperoleh batas bawah.
    Batas atas oleh $\rank T$ mengikuti rumus, dan batas oleh
    $\rank S$ mengikuti $\operatorname{im}(ST)\subseteq\operatorname{im}S$.
\end{proof}

Jika $T$ bijektif, inversnya linear: untuk $w_j=T(v_j)$,
persamaan $T(av_1+bv_2)=aw_1+bw_2$ memberi
$T^{-1}(aw_1+bw_2)=av_1+bv_2$.
Transformasi linear bijektif disebut \emph{isomorfisme}.
Dua ruang berdimensi hingga isomorfik jika dan hanya jika
dimensinya sama: kirim basis ruang pertama ke basis ruang kedua.

\begin{exam}[Invers sebagai jumlah hingga turunan]
    Buktikan bahwa $I-D$ invertibel pada $\mathcal P_{\le n}(\F)$,
    dengan $D(p)=p'$, dan selesaikan $p-p'=q$ di ruang tersebut.
\end{exam}
\begin{solution}
    Karena $D^{n+1}=0$, identitas deret geometri operator memberi
    \[
        (I-D)(I+D+\cdots+D^n)=I-D^{n+1}=I.
    \]
    Perkalian dalam urutan terbalik juga bernilai $I$.
    Jadi $(I-D)^{-1}=I+D+\cdots+D^n$ dan satu-satunya solusi
    adalah $p=q+q'+\cdots+q^{(n)}$.
    Secara umum, jika $N^k=0$, maka
    $(I-N)^{-1}=I+N+\cdots+N^{k-1}$.
\end{solution}

Pada ruang berdimensi hingga yang sama, $ST=I$ mengakibatkan
$TS=I$: persamaan pertama membuat $T$ injektif, sehingga
$T$ bijektif dan $S=T^{-1}$.
Untuk ruang tak hingga, injektif dan surjektif tidak ekuivalen.
Pada $\F[t]$, operator turunan surjektif tetapi tidak injektif,
sedangkan perkalian oleh $t$ injektif tetapi tidak surjektif.

\subsection{Ruang semua transformasi linear}
Himpunan $\mathcal L(V,W)$ menjadi ruang linear dengan operasi
\[
    (S+T)(v)=S(v)+T(v),\qquad (aT)(v)=aT(v).
\]
Jika $\dim V=n$ dan $\dim W=m$, maka
$\dim\mathcal L(V,W)=mn$.
Untuk membuktikannya, pilih basis $(b_j)$ dan $(c_i)$, lalu
definisikan $E_{ij}^{\mathcal L}$ melalui
\[
    E_{ij}^{\mathcal L}(b_k)=\begin{cases}c_i,&k=j,\\0,&k\ne j.\end{cases}
\]
Setiap $T$ mempunyai ekspansi unik
$T=\sum_{i,j}a_{ij}E_{ij}^{\mathcal L}$, dengan $a_{ij}$ entri
matriks representasinya. Ini memberi basis yang beranggotakan $mn$.
Tidak diperlukan hubungan khusus antara vektor basis domain
dan vektor basis kodomain.

\section{Proyeksi dan Subruang Invarian}

\begin{definisi}[][Proyeksi dan subruang invarian][trans:proyeksi]
    Operator $P:V\to V$ disebut \emph{proyeksi} jika $P^2=P$.
    Subruang $U\subseteq V$ disebut \emph{invarian} terhadap
    operator $T$ jika $T(U)\subseteq U$.
\end{definisi}

Proyeksi di sini tidak memerlukan hasil kali dalam. Istilah
proyeksi ortogonal memuat syarat tambahan yang akan bergantung
pada struktur hasil kali dalam.

\begin{teorema}[][Proyeksi dan dekomposisi ruang][trans:dekomposisi-proyeksi]
    Jika $P^2=P$, maka
    \[
        V=\operatorname{im}P\oplus\ker P,\qquad
        \operatorname{im}(I-P)=\ker P.
    \]
    Sebaliknya, jika $V=U\oplus W$, terdapat tepat satu
    proyeksi dengan citra $U$ dan kernel $W$, yaitu $P(u+w)=u$.
\end{teorema}
\begin{proof}
    Setiap $v$ dapat ditulis sebagai $Pv+(v-Pv)$.
    Suku pertama berada di citra, sedangkan suku kedua berada
    di kernel karena $P(v-Pv)=Pv-P^2v=0$.
    Jika $x=Py\in\ker P$, maka $0=Px=P^2y=Py=x$;
    jadi jumlahnya langsung.
    Selain itu, setiap $(I-P)v$ berada di kernel, dan untuk
    $x\in\ker P$ berlaku $x=(I-P)x$.
    Arah sebaliknya mengikuti keunikan dekomposisi $u+w$.
    Proyeksi dengan citra $U$ bertindak sebagai identitas pada $U$
    dan bernilai nol pada $W$, sehingga rumusnya juga unik.
\end{proof}

Jika $\dim V=n$ dan $\rank P=r$, gabungan basis citra dan
kernel memberi matriks $\diag(I_r,0)$.
Maka $\tr P=r$ dan $\rank(I-P)=n-r$.
Untuk operator umum, kernel dan citra memang invarian terhadap
operator itu sendiri, tetapi tidak harus saling komplemen.

\begin{exam}[Memisahkan bagian simetris suatu matriks]
    Pada $\Mat_n(\F)$, tinjau $P(X)=(X+X^T)/2$.
    Tentukan citra, kernel, dan rank dari $P$.
\end{exam}
\begin{solution}
    Transpose linear atas $\F$, sehingga $P$ linear.
    Matriks $P(X)$ simetris, dan untuk matriks simetris $S$
    berlaku $P(S)=S$. Jadi $P^2=P$ dan citranya adalah ruang
    matriks simetris. Kernelnya ruang matriks antisimetris,
    yaitu $X^T=-X$. Dengan menghitung entri bebas,
    \[
        \rank P=\frac{n(n+1)}2,\qquad
        \nul P=\frac{n(n-1)}2.
    \]
    Dekomposisi yang dihasilkan adalah
    $X=(X+X^T)/2+(X-X^T)/2$.
    Jika transpose diganti transpose konjugat, pemetaannya
    linear atas $\R$, tetapi umumnya tidak linear atas $\C$.
\end{solution}

\subsection{Basis yang disesuaikan dengan subruang}
Jika $U$ invarian terhadap $T$, perluas basis $U$ menjadi basis
$V$. Dalam basis ini, matriks $T$ berbentuk
\[
    \begin{pmatrix}A&B\\0&C\end{pmatrix}.
\]
Blok nol muncul karena citra anggota basis $U$ tidak mempunyai
komponen di luar $U$.
Jika $V=U\oplus W$ dan kedua subruang invarian, matriksnya
berbentuk blok diagonal $\diag(A,C)$.
Komplemen suatu subruang selalu ada dalam dimensi hingga,
tetapi komplemen yang invarian belum tentu ada.

Jika $ST=TS$, maka $S$ mempertahankan $\ker T$ dan
$\operatorname{im}T$. Untuk $v\in\ker T$, berlaku
$T(Sv)=S(Tv)=0$; untuk $v=Tu$, berlaku $Sv=T(Su)$.
Argumen yang sama menunjukkan bahwa setiap ruang eigen
$\ker(T-\lambda I)$ invarian terhadap $S$.

\section{Operator Nilpoten dan Stabilisasi Kernel}

Operator $N$ disebut nilpoten jika $N^k=0$ untuk suatu $k\ge1$.
Rantai citra suatu vektor sering menghasilkan daftar bebas linear
tanpa menghitung determinan.

\begin{teorema}[][Rantai nilpoten][trans:rantai]
    Jika $N^kv=0$ dan $N^{k-1}v\ne0$, maka
    \[
        v,Nv,\ldots,N^{k-1}v
    \]
    bebas linear. Khususnya, indeks nilpotensi operator pada
    ruang berdimensi $n\ge1$ paling besar $n$.
\end{teorema}
\begin{proof}
    Andaikan $\sum_{j=0}^{k-1}a_jN^jv=0$ mempunyai koefisien
    tidak nol, dan pilih indeks terkecil $j_0$ dengan $a_{j_0}\ne0$.
    Terapkan $N^{k-1-j_0}$ pada hubungan itu.
    Suku dengan indeks lebih besar lenyap, sehingga
    $a_{j_0}N^{k-1}v=0$, kontradiksi.
    Jika $k$ adalah indeks nilpotensi, ada $v$ dengan
    $N^{k-1}v\ne0$, sehingga $k\le n$.
\end{proof}

\begin{exam}[Basis dari satu rantai turunan]
    Pada $\mathcal P_{\le n}(\F)$, ambil $v=t^n$ dan $D(p)=p'$.
    Tunjukkan bahwa $v,Dv,\ldots,D^nv$ basis dan tentukan
    $\nul D^j$ untuk $0\le j\le n+1$.
\end{exam}
\begin{solution}
    Karena $D^nv=n!\ne0$ dan $D^{n+1}v=0$, teorema rantai
    memberi $n+1$ vektor bebas, sehingga semuanya membentuk basis.
    Untuk $1\le j\le n+1$, kernel $D^j$ terdiri atas polinom
    berderajat paling besar $j-1$, maka $\nul D^j=j$.
    Untuk $j=0$, $D^0=I$ dan nulitasnya nol.
\end{solution}

Untuk operator umum $T:V\to V$ dengan $\dim V=n$, selalu ada
rantai inklusi
\[
    \ker T^0\subseteq\ker T\subseteq\ker T^2\subseteq\cdots,
    \qquad
    V\supseteq\operatorname{im}T\supseteq\operatorname{im}T^2\supseteq\cdots.
\]
Rantai kernel tidak dapat membesar secara ketat lebih dari $n$ kali.
Jika $\ker T^k=\ker T^{k+1}$, maka kesamaan berlaku untuk semua
pangkat selanjutnya. Memang, jika $T^{k+2}v=0$, maka
$Tv\in\ker T^{k+1}=\ker T^k$, sehingga $T^{k+1}v=0$.
Ulangi argumen ini secara induktif.

\begin{teorema}[][Dekomposisi setelah kernel stabil][trans:stabil]
    Misalkan $T$ operator pada $V$ dengan $\dim V=n<\infty$.
    Jika $\ker T^k=\ker T^{k+1}$ untuk suatu $k\ge1$, maka
    \[
        V=\ker T^k\oplus\operatorname{im}T^k.
    \]
    Kedua subruang invarian; pembatasan $T$ pada yang pertama
    nilpoten, dan pada yang kedua invertibel.
    Untuk $\dim V=n\ge1$, pilihan $k=n$ selalu dapat digunakan.
\end{teorema}
\begin{proof}
    Stabilisasi memberi $\ker T^{2k}=\ker T^k$.
    Jika $x=T^ky\in\ker T^k$, maka $T^{2k}y=0$, sehingga
    $y\in\ker T^k$ dan $x=0$.
    Jadi irisan kernel dan citra $T^k$ hanya nol; rank--nulitas
    menyatakan jumlah dimensinya $n$, sehingga jumlahnya seluruh $V$.
    Keduanya invarian karena $T$ komutatif dengan pangkatnya.
    Pada kernel, pangkat ke-$k$ bernilai nol.
    Pada citra, jika $Tx=0$, maka $x\in\ker T^k$ juga,
    sehingga $x=0$. Pembatasan ini injektif pada ruang berdimensi
    hingga, maka invertibel.
    Rantai kernel telah stabil paling lambat pada pangkat $n$.
\end{proof}

Kasus $k=1$ memberi kriteria yang sering berguna:
\[
    V=\ker T\oplus\operatorname{im}T
    \iff\rank T^2=\rank T.
\]
Satu arah mengikuti teorema dan rank--nulitas. Untuk arah lainnya,
jika jumlahnya langsung dan $T^2v=0$, maka
$Tv\in\ker T\cap\operatorname{im}T=\{0\}$, sehingga
$\ker T^2=\ker T$.
Sifat ini tidak memaksa $T$ idempoten: $T=2I$ memenuhi
dekomposisi tersebut, tetapi $T^2\ne T$ pada ruang tak nol.

\section{Latihan Pilihan}

\begin{latihan}\label{trans:lat-data}
    Tentukan semua transformasi linear
    $T:\mathcal P_{\le2}(\R)\to\R^2$ yang memenuhi
    \[
        T(1)=(1,0),\quad T(t)=(0,1),\quad T(1+t)=(a,b).
    \]
    Untuk parameter yang memungkinkan, apakah $T$ unik?
    Tentukan rank dan nulitas setiap transformasi yang mungkin.
\end{latihan}

\begin{latihan}\label{trans:lat-jet}
    Definisikan $T:\mathcal P_{\le4}(\R)\to\R^3$ melalui
    $T(p)=(p(0),p'(0),p(1))^T$.
    Tentukan matriks terhadap basis standar, basis kernel,
    serta rank. Carilah semua $p$ dengan $T(p)=(\alpha,\beta,\gamma)^T$.
\end{latihan}

\begin{latihan}\label{trans:lat-basis}
    Pada $\R^2$, definisikan $T(x,y)=(2x+y,x+2y)$.
    Gunakan basis $\mathcal C=((1,1)^T,(1,-1)^T)$ untuk
    menentukan $[T]_{\mathcal C}$ dan rumus $T^k(x,y)$
    untuk setiap bilangan bulat $k$, termasuk $k<0$.
\end{latihan}

\begin{latihan}\label{trans:lat-euler}
    Untuk $a\in\F$, tinjau operator
    $T_a(p)=tp'-ap$ pada $\mathcal P_{\le n}(\F)$.
    Tentukan kernel, citra, dan syarat invertibilitasnya.
    Jika invertibel, tuliskan $T_a^{-1}$ secara eksplisit.
\end{latihan}

\begin{latihan}\label{trans:lat-komutator}
    Ambil $A=\diag(1,1,2)$ dan definisikan
    $T:\Mat_3(\F)\to\Mat_3(\F)$ melalui $T(X)=AX-XA$.
    Tentukan basis kernel dan citra serta rank dan nulitasnya.
    Generalisasikan hasilnya untuk matriks diagonal dengan
    nilai diagonal berbeda yang bermultiplisitas $m_1,\ldots,m_s$.
\end{latihan}

\begin{latihan}\label{trans:lat-kuadrat-nol}
    Misalkan $T^2=0$ pada ruang berdimensi $n$.
    Buktikan bahwa $\rank T\le\lfloor n/2\rfloor$.
    Tunjukkan bahwa setiap rank bulat antara nol dan batas itu
    dapat dicapai. Kapan $\operatorname{im}T=\ker T$?
\end{latihan}

\begin{latihan}\label{trans:lat-dua-proyeksi}
    Untuk proyeksi $P,Q$ pada ruang yang sama, buktikan bahwa
    $P+Q$ proyeksi jika dan hanya jika $PQ=QP=0$.
    Dalam keadaan itu, buktikan
    \[
        \operatorname{im}(P+Q)=\operatorname{im}P\oplus\operatorname{im}Q.
    \]
\end{latihan}

\begin{latihan}\label{trans:lat-invers-sepihak}
    Misalkan $T:V\to W$ dan $S:W\to V$ linear dengan
    $ST=I_V$, serta kedua ruang berdimensi hingga.
    Buktikan bahwa $TS$ proyeksi pada $W$ dan
    \[
        W=\operatorname{im}T\oplus\ker S.
    \]
    Tentukan kapan $TS=I_W$.
\end{latihan}

\begin{latihan}\label{trans:lat-irisan}
    Untuk transformasi $T:V\to W$ dan subruang $U_1,U_2\subseteq V$,
    buktikan
    \[
        T(U_1+U_2)=T(U_1)+T(U_2),\qquad
        T(U_1\cap U_2)\subseteq T(U_1)\cap T(U_2).
    \]
    Berikan contoh inklusi kedua yang sejati, dan buktikan bahwa
    kesamaan berlaku jika $T$ injektif.
\end{latihan}

\begin{latihan}\label{trans:lat-faktorisasi}
    Misalkan $A:V\to W$ dan $B:V\to Z$ linear, dengan semua
    ruang berdimensi hingga. Buktikan bahwa
    \[
        \ker A\subseteq\ker B
        \iff B=CA\text{ untuk suatu transformasi linear }C:W\to Z.
    \]
    Tunjukkan bahwa $C$ unik jika $A$ surjektif.
\end{latihan}

\begin{latihan}\label{trans:lat-komplemen-invarian}
    Pada $\R^2$, definisikan $N(x,y)=(y,0)$ dan
    $U=\operatorname{span}\{e_1\}$.
    Buktikan bahwa $U$ invarian, tetapi tidak mempunyai
    komplemen invarian terhadap $N$.
\end{latihan}

\begin{latihan}\label{trans:lat-komutator-identitas}
    Pada $\F[t]$, definisikan $D(p)=p'$ dan $M(p)=tp$.
    Buktikan $DM-MD=I$ dan, untuk $k\ge1$,
    \[
        D^kM-MD^k=kD^{k-1}.
    \]
    Buktikan pula bahwa pada ruang tak nol berdimensi hingga
    atas $\R$ atau $\C$, tidak ada operator $S,T$ dengan $ST-TS=I$.
\end{latihan}

\ifpembahasan
\section{Pembahasan Latihan}

\subsection*{Latihan \ref{trans:lat-data}: nilai pada daftar bergantung}
Linearitas memaksa $T(1+t)=T(1)+T(t)=(1,1)$,
sehingga transformasi ada tepat ketika $(a,b)=(1,1)$.
Nilai $T(t^2)=(c,d)$ dapat dipilih sembarang.
Seluruh transformasi yang memenuhi adalah
\[
    T(x+yt+zt^2)=(x+cz,y+dz),\qquad c,d\in\R.
\]
Jadi transformasi tidak unik. Citra memuat $(1,0)$ dan $(0,1)$,
sehingga rank selalu dua dan nulitas satu.
Kernelnya dibangun oleh $t^2-c-dt$.

\subsection*{Latihan \ref{trans:lat-jet}: nilai dan turunan sebagai kendala}
Matriksnya terhadap $(1,t,t^2,t^3,t^4)$ dan basis standar $\R^3$ adalah
\[
    [T]=\begin{pmatrix}1&0&0&0&0\\0&1&0&0&0\\1&1&1&1&1\end{pmatrix}.
\]
Kernel terdiri atas kelipatan $t^2(t-1)$ yang berderajat paling
besar empat. Basisnya $(t^2(t-1),t^3(t-1))$, sehingga nulitas
dua dan rank tiga. Solusi khusus untuk target yang diberikan adalah
\[
    p_0(t)=\alpha(1-t^2)+\beta(t-t^2)+\gamma t^2.
\]
Memeriksa nilai dan turunannya menunjukkan $T(p_0)=(\alpha,\beta,\gamma)^T$.
Semua solusi diperoleh dengan menambahkan kernel:
$p=p_0+t^2(t-1)(c+dt)$, dengan $c,d\in\R$.

\subsection*{Latihan \ref{trans:lat-basis}: pangkat dalam basis yang tepat}
Vektor $c_1=(1,1)^T$ dan $c_2=(1,-1)^T$ memenuhi
$T(c_1)=3c_1$ dan $T(c_2)=c_2$. Jadi $[T]_{\mathcal C}=\diag(3,1)$.
Karena kedua entri diagonal tak nol, $T$ invertibel dan
pangkat bilangan bulat terdefinisi. Ekspansi
\[
    (x,y)^T=\frac{x+y}{2}c_1+\frac{x-y}{2}c_2
\]
memberi
\[
    T^k(x,y)=\left(
      \frac{3^k(x+y)+x-y}{2},\
      \frac{3^k(x+y)-x+y}{2}\right).
\]

\subsection*{Latihan \ref{trans:lat-euler}: bekerja pada setiap monomial}
Karena $T_a(t^j)=(j-a)t^j$, matriksnya terhadap basis monomial
adalah $\diag(-a,1-a,\ldots,n-a)$.
Jika $a\notin\{0,1,\ldots,n\}$, semua entri tak nol,
kernel hanya nol, citra seluruh ruang, dan
\[
    T_a^{-1}\left(\sum_{j=0}^n b_jt^j\right)
      =\sum_{j=0}^n\frac{b_j}{j-a}t^j.
\]
Jika $a=r\in\{0,\ldots,n\}$, kernel dibangun oleh $t^r$ dan
citra terdiri atas polinom yang koefisien $t^r$-nya nol.
Ranknya $n$, nulitasnya satu, dan operator tidak invertibel.

\subsection*{Latihan \ref{trans:lat-komutator}: matriks satuan sebagai basis}
Untuk $A=\diag(\lambda_1,\ldots,\lambda_n)$, berlaku
\[
    T(E_{ij})=(\lambda_i-\lambda_j)E_{ij}.
\]
Pada kasus $\diag(1,1,2)$, basis kernel adalah
$(E_{11},E_{12},E_{21},E_{22},E_{33})$, sedangkan basis citra
adalah $(E_{13},E_{23},E_{31},E_{32})$.
Jadi nulitas lima dan rank empat.
Secara umum, kernel dibangun oleh $E_{ij}$ dengan
$\lambda_i=\lambda_j$ dan citra oleh sisanya. Kelompok nilai
diagonal bermultiplisitas $m_\ell$ menyumbang $m_\ell^2$ anggota
basis kernel. Maka
\[
    \nul T=\sum_{\ell=1}^s m_\ell^2,\qquad
    \rank T=n^2-\sum_{\ell=1}^s m_\ell^2.
\]

\subsection*{Latihan \ref{trans:lat-kuadrat-nol}: citra termuat dalam kernel}
Dari $T^2=0$, diperoleh $\operatorname{im}T\subseteq\ker T$.
Jika $r=\rank T$, maka $r\le n-r$, sehingga $r\le\lfloor n/2\rfloor$.
Untuk setiap $r$ dalam rentang itu, pilih basis $(e_1,\ldots,e_n)$
dan tetapkan $T(e_{2j})=e_{2j-1}$ untuk $1\le j\le r$,
serta $T(e_k)=0$ untuk semua indeks lainnya.
Operator ini berpangkat dua nol dan citranya dibangun oleh
$e_1,e_3,\ldots,e_{2r-1}$, sehingga ranknya $r$.
Inklusi citra dalam kernel menjadi kesamaan tepat ketika
dimensinya sama, yaitu $n=2r$.

\subsection*{Latihan \ref{trans:lat-dua-proyeksi}: dua perkalian harus lenyap}
Jika $P+Q$ idempoten, ekspansi kuadrat memberi $PQ+QP=0$.
Kalikan di kiri dan di kanan oleh $P$, secara terpisah:
\[
    PQ+PQP=0,\qquad PQP+QP=0.
\]
Mengurangkan keduanya memberi $PQ=QP$.
Karena $PQ+QP=0$ dan medan berkarakteristik nol, diperoleh
$PQ=QP=0$. Arah sebaliknya langsung dari ekspansi kuadrat.
Jika $x\in\operatorname{im}P\cap\operatorname{im}Q$, maka
$Px=x$, tetapi $x=Qy$ memberi $Px=PQy=0$; jadi $x=0$.
Untuk $u=Pa$ dan $v=Qb$ berlaku $(P+Q)(u+v)=u+v$.
Maka setiap anggota jumlah citra berada di citra $P+Q$;
inklusi sebaliknya langsung dari rumus $P+Q$.

\subsection*{Latihan \ref{trans:lat-invers-sepihak}: proyeksi pada ruang yang lebih besar}
Persamaan $ST=I_V$ membuat $T$ injektif dan $S$ surjektif.
Selain itu, $(TS)^2=T(ST)S=TS$, sehingga $TS$ proyeksi.
Citranya sama dengan citra $T$: satu inklusi jelas, dan
untuk $Tv$ berlaku $TS(Tv)=Tv$.
Kernelnya sama dengan kernel $S$, sebab $TSw=0$ dan injektivitas
$T$ mengakibatkan $Sw=0$.
Teorema proyeksi memberi $W=\operatorname{im}T\oplus\ker S$.
Selanjutnya, $TS=I_W$ tepat ketika $T$ surjektif, atau ekuivalen
dengan $S$ injektif. Dalam dimensi hingga dan dengan $ST=I_V$,
ini terjadi tepat ketika $\dim V=\dim W$.

\subsection*{Latihan \ref{trans:lat-irisan}: prapeta yang berbeda}
Linearitas memberi $T(u_1+u_2)=T(u_1)+T(u_2)$, sehingga
kesamaan jumlah berlaku. Jika $v\in U_1\cap U_2$, citranya
berada dalam kedua citra, membuktikan inklusi irisan.
Untuk inklusi sejati, ambil $T:\R^2\to\R$,
$T(x,y)=x+y$, serta $U_1=\operatorname{span}\{e_1\}$ dan
$U_2=\operatorname{span}\{e_2\}$.
Irisan sumber hanya nol, tetapi kedua citra sama dengan $\R$.
Jika $T$ injektif dan $w=T(u_1)=T(u_2)$ dengan $u_j\in U_j$,
maka $u_1=u_2\in U_1\cap U_2$, sehingga inklusi balik berlaku.

\subsection*{Latihan \ref{trans:lat-faktorisasi}: definisikan dahulu pada citra}
Jika $B=CA$, setiap $v\in\ker A$ memenuhi $Bv=C(Av)=0$.
Sebaliknya, andaikan $\ker A\subseteq\ker B$.
Definisikan $C_0:\operatorname{im}A\to Z$ melalui
$C_0(Av)=Bv$.
Jika $Av=Av'$, maka $v-v'\in\ker A\subseteq\ker B$,
sehingga $Bv=Bv'$; jadi definisinya sah.
Linearitas mengikuti dari $A$ dan $B$.
Perluas basis citra $A$ menjadi basis $W$, lalu tentukan
nilai $C$ pada vektor tambahan secara sembarang, misalnya nol.
Teorema penentuan pada basis memberi perluasan linear dengan $B=CA$.
Jika $A$ surjektif, semua anggota $W$ berbentuk $Av$ dan
nilainya sudah dipaksa oleh $C(Av)=Bv$, sehingga $C$ unik.

\subsection*{Latihan \ref{trans:lat-komplemen-invarian}: komplemen tidak selalu terpelihara}
Karena $N(e_1)=0$, subruang $U$ invarian.
Setiap komplemennya berupa garis yang dibangun oleh
$w=(a,b)$ dengan $b\ne0$.
Jika garis ini invarian, harus ada $\lambda\in\R$ dengan
$Nw=(b,0)=\lambda(a,b)$.
Koordinat kedua memaksa $\lambda=0$, sedangkan koordinat
pertama kemudian memaksa $b=0$, kontradiksi.

\subsection*{Latihan \ref{trans:lat-komutator-identitas}: perbedaan dimensi hingga dan tak hingga}
Aturan turunan hasil kali memberi
$(DM-MD)p=(tp)'-tp'=p$, sehingga $DM-MD=I$.
Untuk $k\ge1$, gunakan identitas teleskopik
\[
    D^kM-MD^k
      =\sum_{j=0}^{k-1}D^j(DM-MD)D^{k-1-j}
      =kD^{k-1}.
\]
Pada ruang berdimensi $n\ge1$, representasi matriks akan
memberikan $AB-BA=I_n$ jika operator demikian ada.
Mengambil trace memberi $0=\tr(AB)-\tr(BA)=n$, kontradiksi
di atas $\R$ atau $\C$.
Tidak ada pertentangan dengan contoh polinom: $\F[t]$
berdimensi tak hingga, dan perkalian oleh $t$ tidak mempertahankan
subruang $\mathcal P_{\le n}(\F)$.
\fi
